\section{Owned choice and rational outcomes}
\label{sec:rational}

Fix a finite set $I$ of $K\geq3$ roles. A role values its own win at one
and every other outcome at zero. Values are nonempty winner sets
$\Pow(I)$. A singleton is \emph{resolved}; a set with at least two members
is \emph{ambiguous}. The full term language is
\[
 t ::= \at i\mid x_j\mid t\ch r t\mid\pi(t),
 \qquad i,r\in I,\quad \pi\in\Sym(I).
\]
The constant $\at i$ denotes $\{i\}$. The constructor $\ch r$ is
owned choice, and $\pi(t)$ internally permutes the component's roles.
An input variable is evaluated in its original environment before that
permutation is applied. Substitution at repeated occurrences creates
independent component copies.
Leaves carry only a winner label, with no strategic terminal returns.
The dependent interface in \cref{app:dependent} also permits return
records; a simultaneous decision at such a boundary is outside this
perfect-information winner-set recurrence.

We distinguish three signatures throughout:
\[
 \begin{array}{c|l}
 \mathbf A_I & \star_r\ (r\in I),\ \{i\}\ (i\in I)\\
 \mathbf A_I^{\rm sym} & \mathbf A_I\text{ and all internal role permutations}\\
 \mathbf C_I & \star_r\ (r\in I)\text{ alone, with variables and no constants}.
 \end{array}
\]
All have carrier $\Pow(I)$. External renaming of the roles is a symmetry
in every case; an internal permutation is an extra term constructor only
in the full signature. The core $\mathbf C_I$ will be characterized
intrinsically by idempotency in \cref{sec:intrinsic}.

\subsection{From full strategies to a finite algebra}

A full profile gives an action at \emph{every} internal node, including
nodes avoided by that profile. At a node it therefore consists of a root
action and a full profile for each child. A unilateral deviation by $i$
may replace all actions owned by $i$, in both child trees. It leaves every
other player's actions fixed. A profile is Nash if no such deviation
improves the deviator's utility. It is SPE if its restriction to every
subtree is Nash.

For a fixed profile, let $a,b$ be the winners prescribed in its two child
subtrees. Assuming both child profiles are SPE, the parent profile is SPE
exactly when it chooses $r$ whenever $a=r$ or $b=r$. Otherwise either
child is optimal for $r$. This characterization follows from whole-player
deviations: a deviation that wins in a child would already improve that
child's profile. It is not an assumption that checking only a current
action is generally sufficient in games of partial information.

For nonempty subsets $S,T\subseteq I$, put
\begin{equation}\label{game:eq:own}
 S\ow r T =
 \begin{cases}
 \{r\},&S=\{r\}\text{ or }T=\{r\},\\
 S\cup T,&\text{otherwise}.
 \end{cases}
\end{equation}

\begin{theorem}[Exact public-tree interpretation]\label{game:thm:adequacy}
Every closed tree has a full pure SPE profile, and
\[
 \B(\at i)=\{i\},\qquad
 \B(G\ch r H)=\B(G)\ow r\B(H).
\]
Every element of the right-hand side has a full SPE witness on the left.
\end{theorem}
\begin{proof}
Existence follows inductively by choosing an owner-winning child whenever
one is available in the chosen child profiles. If one child has only
$r$-winning SPE profiles, every parent SPE must win for $r$, and one
exists. Otherwise each child has an SPE profile losing for $r$. For any
SPE winner $i$ in the left child, combine its witness with a losing-for-$r$
SPE in the right child and choose left. This is locally optimal, whether
or not $i=r$. The symmetric construction covers the right child.
Conversely, a parent SPE selects the winner of a child SPE. These facts
give both inclusions in~\eqref{game:eq:own}.
\end{proof}

This is a backward-induction result for a restricted utility model. Its
role here is to justify the finite algebra $\Pow(I)$, not to claim new
game theory. \texttt{Basic.lean} defines full profiles, whole-player
deviations, Nash equilibrium, and SPE separately, then proves the theorem
as \texttt{winners\_node}.

Each $\star_r$ is associative, commutative, and idempotent. For
associativity, a three-input product is $\{r\}$ if any input is $\{r\}$,
and is their union otherwise. Nonemptiness ensures that a union of inputs
different from $\{r\}$ cannot become $\{r\}$. Role permutations transport
these operations and the leaf constants.

For two roles, the ground-generated carrier consists only of the two
singletons and recovers Boolean conjunction/disjunction after choosing a
role as truth. With three roles there is genuine ambiguity, such as
$\{0,1\}$, and the next construction generates every nonempty subset.

\subsection{Neutral union and its minimum occurrence cost}

\begin{theorem}[Neutral union]\label{game:thm:union}
For pairwise distinct $r,s,t$,
\begin{equation}\label{game:eq:union}
 J(x,y)=(x\ow r y)\ow t(x\ow s y)=x\cup y.
\end{equation}
Consequently, if $K\geq3$, every nonempty subset of $I$ is the
denotation of a closed owned-choice tree.
\end{theorem}
\begin{proof}
If neither inner operation collapses, both return $x\cup y$. If just one
collapses, its owner singleton belongs to the other result, and the outer
owner is different. If both collapse, the inputs must be the two distinct
owner singletons; the outer operation returns their union. Folding $J$
over the leaf singletons realizes every nonempty subset.
\end{proof}

The definition uses three owned-choice nodes and two occurrences of each
input. Although it names three owners, its denotation is independent of
the distinct triple chosen. It is an equality of equilibrium-winner
observations, not an isomorphism of trees or strategy spaces.

\begin{theorem}[Occurrence lower bound]\label{game:thm:lower}
Suppose $K\geq3$. In the syntax of owned binary choices and arbitrary
nonempty ground constants, any term that denotes $x\cup y$ for every
$x,y\in\Pow(I)$ contains at least two occurrences of $x$ and at least two
occurrences of $y$. Hence three binary nodes are optimal.
\end{theorem}
\begin{proof}
Zero occurrences cannot express dependence on that input. Suppose $x$
occurs once. Read the unique path from its occurrence to the root, using
commutativity to put the current value first. It has the form
\[
 X\longmapsto X\ow{r_1}g_1(Y)
   \longmapsto\cdots\longmapsto (\cdots)\ow{r_n}g_n(Y),
\]
where each $g_j(Y)$ is nonempty and independent of $X$. We allow arbitrary
such functions, which only strengthens the obstruction.

For fixed $Y$, if some $g_j(Y)=\{r_j\}$, that step erases all dependence
on $X$, and later steps cannot recover it. For $Y=\{r\}$ this is
incompatible with $X\cup\{r\}$, which varies with $X$. Therefore no such
side value equals its owner singleton. Every step then preserves all
members of its current input: if the input is the owner singleton it
remains so, and otherwise the operation takes a union.

If the path is empty, its value is $X$, not union. Otherwise let $r$ be
its innermost owner and set $Y=\{r\}$. Its first side value differs from
$\{r\}$ and is nonempty, so contains some $k\ne r$. Choose
$j\notin\{r,k\}$. Starting with $X=\{j\}$, the first step introduces
$k$. No later step can remove it. The required result $\{j,r\}$ excludes
$k$, a contradiction. Interchange the variables for the other bound.
\end{proof}

The path obstruction is \texttt{no\_single\_use\_union} in Lean. Ground
game realization is \texttt{every\_winner\_set\_realized} in
\texttt{Synthesis.lean}. The
syntax-to-path correspondence is the direct decomposition of a term at
its unique input occurrence, as used in the proof.

The cost measured here is \emph{syntactic occurrence}. A play enters just
one of the four component occurrences. The theorem establishes neither
runtime duplication nor a need for a contraction rule in a particular
linear type system: additive branching may use a context in both
exclusive branches.


\subsection{A distributivity obstruction}
For three roles, indifference gives additional values and invalidates
classical distributivity. In particular,
\[
 B((\langle0\rangle\oplus_0\langle1\rangle)
                  \oplus_1\langle2\rangle)=\{0,2\},
\]
whereas
\[
 B((\langle0\rangle\oplus_1\langle2\rangle)
           \oplus_0(\langle1\rangle\oplus_1\langle2\rangle))
       =\{0,1,2\}.
\]
In the first tree, role $0$ must choose winner $0$ in its subtree.
The root owner $1$ then chooses between prescribed winners $0$ and $2$,
and is indifferent between them. For the distributed tree, the right
subtree's owner $1$ must choose winner $1$, while the left subtree's
owner $1$ is indifferent between winners $0$ and $2$. Its full SPE
profiles are therefore exactly the following three choices of child
outcomes and root action:
\begin{center}
\begin{tabular}{cccc}
\toprule
Left child winner & Right child winner & Root action (owner $0$) & Winner\\
\midrule
$0$ & $1$ & left & $0$\\
$2$ & $1$ & left & $2$\\
$2$ & $1$ & right & $1$\\
\bottomrule
\end{tabular}
\end{center}
The last row prescribes winner $2$ in the unvisited left subtree.
That continuation is an equilibrium because its owner $1$ loses either
way; it makes root owner $0$ indifferent between its two losing choices.
This is how the repeated branch enables the additional winner $1$.
The calculation concerns full subgame-perfect profiles, including their
unvisited continuations.

\begin{corollary}[No Boolean meet/join product representation]
\label{cor:no-boolean-product}
For $K\ge3$, this generated algebra has no injective homomorphism into a
product of Boolean algebras in which each owned operation acts in each
coordinate as either meet or join.
\end{corollary}
\begin{proof}
Any pair of such coordinate operations satisfies
$(a\star_r b)\star_s c
 =(a\star_s c)\star_r(b\star_s c)$, whether the coordinate operations
coincide or differ. Products and subalgebras preserve this identity.
The displayed three-role calculation refutes it in the generated algebra.
An injective homomorphism would preserve and reflect the equality.
\end{proof}

For any number of roles, a quantifier over a finite nonempty action set
is the fold of its owner's associative, commutative, idempotent operation.
It adds no independent algebraic primitive. At $K=2$ the ground-generated
carrier has only the two singletons, and the operations are Boolean
join and meet after selecting a viewpoint. The theorems below concern
$K\geq3$, where all nonempty subsets are generated.
